% =========================================================================
\section{El potencial eléctrico}\label{sec:potencial}
% =========================================================================
Hasta ahora hemos descrito los efectos electrostáticos únicamente mediante el campo vectorial $\vE$. Escribiendo el campo de una
carga puntual de otra forma podremos expresar la misma información mediante un \emph{campo escalar}, lo cual es más
conveniente para muchos propósitos: un número por punto en lugar de tres.

% -------------------------------------------------------------------------
\subsection{Definición del potencial}\label{sec:def-potencial}
% -------------------------------------------------------------------------
Partimos del campo de una carga puntual (\cref{eq:campo-puntual}) y usamos que $\vc{R}/R^{3}=-\nab(1/R)$, que es el caso $n=-1$ del
gradiente de $R^{n}$~\cite{TEMA0} ($\nab R^{n}=nR^{n-1}\uR$):
\begin{equation}\label{eq:campo-grad}
  \vE(\vc{r})=\frac{\qp}{4\pi\eo}\,\frac{\vc{R}}{R^{3}}=-\frac{\qp}{4\pi\eo}\,\nab\Bigl(\frac{1}{R}\Bigr)
  =-\nab\Bigl(\frac{\qp}{4\pi\eo R}\Bigr).
\end{equation}

\begin{definicion}[Potencial electrostático]\label{def:potencial}
Se llama \emph{potencial eléctrico} (o escalar) de una carga puntual $\qp$ al campo escalar
\begin{equation}\label{eq:potencial-puntual}
  \phi(\vc{r})=\frac{1}{4\pi\eo}\,\frac{\qp}{R}+\text{cte},\qquad R=\norm{\vc{r}-\vc{r}\,'}.
\end{equation}
La relación entre el campo vectorial $\vE:\R^{3}\to\R^{3}$ y el campo escalar $\phi:\R^{3}\to\R$ es
\begin{equation}\label{eq:E-grad-phi}
  \boxed{\;\vE(\vc{r})=-\nab\phi(\vc{r})\;}
\end{equation}
\end{definicion}

\begin{observacion}[Propiedades del potencial]\label{obs:potencial}
\begin{enumerate}[label=(\alph*)]
  \item \emph{Significado del gradiente.} $\norm{\nab\phi}$ es la derivada de $\phi$ en la dirección de máxima variación, y $\nab\phi$
        apunta en esa dirección, que coincide con la de $\vE$ en cada punto; pero, por el signo menos, $\vE$ y $\nabla\phi$ tienen
        sentidos opuestos: el potencial \emph{aumenta} al movernos en sentido contrario al campo. Como $\dif\phi=\nab\phi\cdot\dl$,
        resulta $\dif\phi=-\vE\cdot\dl$.
  \item \emph{Unidades.} En el SI, $[\phi]=\mathrm{V}$ (voltio $=\mathrm{N\,m/C}$), y de aquí $[\vE]=\mathrm{V/m}$.
  \item \emph{No está unívocamente definido.} Sumándole cualquier constante se obtiene el mismo campo (el gradiente de una constante
        es nulo). Ocurre como con la altitud de un lugar: siempre se da respecto de una referencia, el nivel del mar. En electrostática se toma
        habitualmente como nivel de referencia (al que se asigna el valor cero) un punto infinitamente alejado de las cargas: si
        $r\to\infty$, $\phi=\lim\bigl(\frac{1}{4\pi\eo}\frac{\qp}{R}+\text{cte}\bigr)=\text{cte}$, luego $\text{cte}=0$. Pero hay
        problemas ideales en los que las cargas llegan al infinito (un hilo o un plano infinitos); en ellos hay que tomar otro punto como
        origen de potencial, e incluso a veces se asigna el cero a la superficie de la Tierra. En cada problema hay que especificarlo, y
        como el potencial no es único conviene hablar de \emph{diferencias de potencial}. A partir de ahora escribimos $\phi$ para el potencial con
        la constante ya elegida.
\end{enumerate}
\end{observacion}

Por el principio de superposición, el potencial de un sistema de cargas es la suma de los potenciales de cada una, y el de una
distribución continua la integral correspondiente:

\begin{proposicion}[Potencial de sistemas y distribuciones de carga]\label{prop:potencial-dist}
\begin{align}
  \phi(\vc{r})&=\frac{1}{4\pi\eo}\sum_{i=1}^{N}\frac{\qp_i}{R_i}+\text{cte}, &&\text{(discreta)}\label{eq:pot-discreto}\\
  \phi(\vc{r})&=\frac{1}{4\pi\eo}\int_{V'}\frac{\rho(\vc{r}\,')\,\dVp}{R}+\text{cte}, &&\text{(volúmica)}\label{eq:pot-vol}\\
  \phi(\vc{r})&=\frac{1}{4\pi\eo}\int_{S'}\frac{\sigma(\vc{r}\,')\,\dSp}{R}+\text{cte}, \quad
  \phi(\vc{r})=\frac{1}{4\pi\eo}\int_{l'}\frac{\lambda(\vc{r}\,')\,\dlp}{R}+\text{cte}. &&\text{(superficial, lineal)}\label{eq:pot-sup-lin}
\end{align}
\end{proposicion}
\begin{proof}
El gradiente es lineal y $\vE=-\nab\phi$ para cada carga puntual, con la misma constante aditiva; la superposición de campos
(\cref{eq:campo-discreto,eq:campo-vol,eq:campo-sup-lin}) es entonces la superposición de potenciales.
\end{proof}

% -------------------------------------------------------------------------
\subsection{Significado físico: el potencial como trabajo por unidad de carga}\label{sec:potencial-signif}
% -------------------------------------------------------------------------
Consideremos una carga puntual $q$ (de prueba, $q\to0$) en equilibrio bajo la acción de la fuerza electrostática $\vF_q=q\vE$ y de
una fuerza \emph{mecánica} externa $\vF$ (\cref{fig:trabajo}a): $\vF_q+\vF=q\vE+\vF=\vc{0}$. Supongamos que la carga se mueve
\emph{infinitamente despacio} (de forma reversible) de un punto $A=\vc{r}_1$ a otro $B=\vc{r}_2$ por una trayectoria dada $\Gamma$. La
velocidad es esencialmente nula en todo momento, también la energía cinética, y la carga está siempre en equilibrio o casi, de modo que la
ecuación anterior es aplicable. Este procedimiento permite calcular el trabajo reversible, en sentido termodinámico, que realiza la fuerza
mecánica contra las del campo, sin efectos disipativos ni de fricción.

% Figura: trabajo de la fuerza externa contra el campo y circulación nula.
\begin{figure}[ht]
\centering
\begin{minipage}[t]{0.52\textwidth}\centering
\begin{tikzpicture}[baseline=(current bounding box.north), x=1cm, y=1cm]
  \coordinate (A) at (0,0); \coordinate (B) at (5.2,2.2);
  \draw[trayecto] (A) .. controls (0.8,2.2) and (2.0,-0.9) .. (3.0,0.9) .. controls (3.6,1.9) and (4.4,1.2) .. (B);
  \draw[vect, cm1, line width=0.9pt] (3.0,0.9) -- ++(-0.15,-0.9) node[below, font=\small, text=cm1] {$\vc{F}_{q}=q\vE$};
  \draw[vectrojo, line width=0.9pt] (3.0,0.9) -- ++(0.15,0.9) node[above, font=\small, text=acento] {$\vc{F}=-q\vE$};
  \draw[-{Stealth[length=4pt]}, azulprof, line width=0.8pt] (3.0,0.9) -- ++(0.55,0.72) node[right, font=\scriptsize] {$\dl$};
  \node[cargapru, minimum size=8pt] at (3.0,0.9) {};
  \fill[azulprof] (A) circle (1.8pt) node[below left, font=\small] {$A$};
  \fill[azulprof] (B) circle (1.8pt) node[above right, font=\small] {$B$};
  \node[font=\small, text=acento, anchor=north west] at (0.7,2.1) {$\Gamma$};
\end{tikzpicture}\\[2pt]{\small\scshape (a)}
\end{minipage}\hfill
\begin{minipage}[t]{0.44\textwidth}\centering
\begin{tikzpicture}[baseline=(current bounding box.north), x=1cm, y=1cm]
  \draw[trayecto, orienta] (0,0) .. controls (0.7,1.8) and (2.3,2.0) .. (3.2,0.2);
  \draw[trayecto, orienta, azulacero] (3.2,0.2) .. controls (2.8,-1.2) and (0.7,-1.5) .. (0,0);
  \fill[azulprof] (0,0) circle (1.8pt) node[left, font=\small] {$1$};
  \fill[azulprof] (3.2,0.2) circle (1.8pt) node[right, font=\small] {$2$};
  \node[font=\small, text=acento] at (1.6,1.75) {$1\to2$};
  \node[font=\small, text=azulacero] at (1.6,-1.2) {$2\to1$};
  \node[font=\small, text=azulprof] at (1.6,0.2) {$C$};
\end{tikzpicture}\\[2pt]{\small\scshape (b)}
\end{minipage}
\caption{(a)~Una fuerza externa $\vF$ lleva, muy despacio, una carga de prueba de $A$ a $B$
contra la fuerza del campo $q\vE$; en cada punto ambas se compensan. (b)~Un circuito cerrado $C$ se
descompone en dos tramos de $1$ a $2$ y de $2$ a $1$; si la integral de línea de $\vE$ no depende del
camino, la circulación a lo largo de $C$ es nula.}
\label{fig:trabajo}
\end{figure}


\begin{proposicion}[Trabajo y diferencia de potencial]\label{prop:trabajo}
El trabajo realizado por la fuerza mecánica $\vF=-q\vE$ para llevar la carga $q$ de $A$ a $B$ por cualquier camino $\Gamma$ es
\begin{equation}\label{eq:trabajo}
  W_{AB}=\int_{A}^{B}\vF\cdot\dl=-q\int_{A}^{B}\vE\cdot\dl=q\int_{A}^{B}\nab\phi\cdot\dl=q\bigl[\phi(B)-\phi(A)\bigr].
\end{equation}
En particular, para mover una carga \emph{positiva unidad} desde el infinito hasta un punto $P$ (con $\phi(\infty)=0$),
\begin{equation}\label{eq:trabajo-unidad}
  W=\int_{\infty}^{P}\vF\cdot\dif\vc{r}=-\int_{\infty}^{P}\vE\cdot\dif\vc{r}=\int_{\infty}^{P}\dif\phi=\phi_P-\phi_\infty=\phi_P.
\end{equation}
\end{proposicion}
\begin{proof}
Se sustituye $\vF=-q\vE$, se usa $\vE=-\nab\phi$ y el teorema fundamental de las integrales de línea (\emph{teorema del gradiente}~\cite{TEMA0}):
$\int_A^B\nab\phi\cdot\dl=\phi(B)-\phi(A)$, que no depende de $\Gamma$.
\end{proof}

El potencial no es, pues, más que el \emph{trabajo por unidad de carga} necesario para desplazar una carga unidad desde el infinito
hasta el punto, \emph{sin importar el camino seguido} (si la fuente es positiva, es trabajo realizado contra el campo). Si el
trabajo suministrado fuese mayor que $q\phi_P$, el exceso de energía se transformaría en energía cinética, el sistema dejaría de estar
en reposo y abandonaríamos la electrostática.

% -------------------------------------------------------------------------
\subsection{El campo electrostático es conservativo e irrotacional}\label{sec:rotacional}
% -------------------------------------------------------------------------
\begin{teorema}[Circulación nula: el campo electrostático es conservativo]\label{teo:conservativo}
Entre dos puntos arbitrarios,
\begin{equation}\label{eq:int-linea}
  -\int_{1}^{2}\vE\cdot\dif\vc{r}=\int_{1}^{2}\nab\phi\cdot\dif\vc{r}=\int_{1}^{2}\dif\phi=\phi(\vc{r}_2)-\phi(\vc{r}_1),
\end{equation}
luego la integral de línea de $\vE$ es independiente del camino seguido y, para cualquier circuito cerrado $C$,
\begin{equation}\label{eq:circulacion}
  \oint_{C}\vE\cdot\dif\vc{r}=\int_{1}^{2}\vE\cdot\dif\vc{r}+\int_{2}^{1}\vE\cdot\dif\vc{r}=0 .
\end{equation}
\end{teorema}
\begin{proof}
La primera cadena de igualdades es el teorema del gradiente. Para la segunda, descomponemos $C$ en dos tramos de $1$ a $2$ y de $2$ a $1$
(\cref{fig:trabajo}b); como $\int_{2}^{1}=-\int_{1}^{2}$ y la integral entre 1 y 2 no depende del camino, la suma es cero.
\end{proof}

\begin{teorema}[El campo electrostático es irrotacional]\label{teo:irrotacional}
En electrostática,
\begin{equation}\label{eq:rotacional}
  \boxed{\;\nab\times\vE=\vc{0}\;}
\end{equation}
\end{teorema}
\begin{proof}
Por el teorema de Stokes~\cite{TEMA0} y \cref{eq:circulacion}, para toda superficie $S$ con borde $C=\partial S$,
\[
  \oint_{C}\vE\cdot\dif\vc{r}=\iint_{S}(\nab\times\vE)\cdot\dS=0 .
\]
El resultado es válido para cualquier trayectoria $C$ y, por tanto, para cualquier superficie $S$; como en la demostración de la ley
de Gauss diferencial, esto obliga a que el integrando $\nab\times\vE$ sea idénticamente nulo.
\end{proof}

\begin{observacion}[Comprobación directa y alcance]\label{obs:rotacional}
La anulación del rotacional se comprueba también directamente: para una distribución volúmica,
$\nab\times\vE=\frac{1}{4\pi\eo}\int_{V'}\rho(\vc{r}\,')\,\nab\times\bigl(\vc{R}/R^{3}\bigr)\dVp=\vc{0}$, pues
$\nab\times(R^{n}\uR)=\vc{0}$ para todo $n$ (caso $n=-2$)~\cite{TEMA0}. \emph{Cuidado}: $\nab\times\vE=\vc{0}$ es un resultado de la
\emph{electrostática}; con campos magnéticos variables en el tiempo (Faraday) dejará de cumplirse. Y a la inversa: si $\vE=-\nab\phi$, entonces
$\nab\times\vE=-\nab\times(\nab\phi)=\vc{0}$ porque el rotacional de un gradiente siempre es nulo (si $\phi$ es suficientemente regular);
un campo con $\nab\times\vE=\vc{0}$ en una región simplemente conexa puede escribirse como el gradiente de un campo escalar~\cite{TEMA0}.
\end{observacion}

Conocidas la divergencia y el rotacional de un campo vectorial en todos los puntos de una región, el campo queda determinado. Es el contenido
del siguiente resultado.

\begin{teorema}[Teorema de Helmholtz~\cite{HELMHOLTZ}]\label{teo:helmholtz}
Sea $\vc{F}$ un campo vectorial con $\nab\cdot\vc{F}=b(\vc{r})$ y $\nab\times\vc{F}=\vc{c}(\vc{r})$ dados en todo el espacio, y tal que $\vc{F}\to\vc{0}$
en el infinito suficientemente deprisa (y las fuentes $b,\vc{c}$ están localizadas). Entonces $\vc{F}$ queda determinado de forma única.
\end{teorema}

En electrostática, $\nab\cdot\vE=\rho/\eo$ y $\nab\times\vE=\vc{0}$ (con las condiciones de regularidad del teorema, que el campo cumple), de modo que el
campo electrostático está \emph{unívocamente determinado} por la distribución de carga. Es la razón por la que la ley de Gauss en forma diferencial y la
irrotacionalidad contienen toda la información del campo, y no solo una parte.

% -------------------------------------------------------------------------
\subsection{Ecuaciones de Poisson y de Laplace}\label{sec:poisson}
% -------------------------------------------------------------------------
¿Qué debemos conocer para calcular $\phi$? Sustituimos $\vE=-\nab\phi$ en las ecuaciones del campo.

\begin{teorema}[Ecuaciones de Poisson y de Laplace]\label{teo:poisson}
El potencial electrostático satisface la \emph{ecuación de Poisson}
\begin{equation}\label{eq:poisson}
  \boxed{\;\Lap\phi=-\frac{\rho}{\eo}\;}
\end{equation}
y, en recintos sin fuentes ($\rho=0$), la \emph{ecuación de Laplace}
\begin{equation}\label{eq:laplace}
  \Lap\phi=0 .
\end{equation}
\end{teorema}
\begin{proof}
Con $\vE=-\nab\phi$ en la ley de Gauss diferencial,
$\nab\cdot\vE=\nab\cdot\bigl(-\nab\phi\bigr)=-\Lap\phi=\rho/\eo$. En una región con $\rho=0$ se obtiene \cref{eq:laplace}.
\end{proof}

Si en lugar de la divergencia sustituimos en el rotacional, $\nab\times\vE=\nab\times(-\nab\phi)=\vc{0}$ se cumple idénticamente: la irrotacionalidad
\emph{no aporta} ninguna ecuación adicional para $\phi$; basta una sola ecuación, la de Poisson. Es una ventaja práctica: para los campos vectoriales
hacen falta dos ecuaciones (divergencia y rotacional), pero para los escalares solo una, y al ser $\phi$ una función escalar resulta más sencillo resolverla y
calcular después $\vE=-\nab\phi$. Para una carga puntual, $\Lap\phi=\frac{\qp}{4\pi\eo}\Lap(1/R)=-\frac{\qp}{\eo}\delta^{3}(\vc{R})$, coherente con
el \cref{teo:delta}.

% -------------------------------------------------------------------------
\subsection{Líneas de campo y superficies equipotenciales}\label{sec:lineas}
% -------------------------------------------------------------------------
El campo eléctrico se suele representar, en cada punto, mediante un vector; globalmente, mediante sus líneas de fuerza.

\begin{definicion}[Líneas de campo y superficies equipotenciales]\label{def:lineas}
Las \emph{líneas de fuerza} (o \emph{de campo}) son las líneas que en cada punto son tangentes al vector $\vE$ en ese punto. Las
\emph{superficies equipotenciales} son aquellas en las que $\phi(\vc{r})$ es constante.
\end{definicion}

\begin{proposicion}[Propiedades de las líneas de campo y de las equipotenciales]\label{prop:lineas}
\begin{enumerate}[label=(\alph*)]
  \item Las líneas de fuerza \emph{no se cortan}: si lo hicieran, $\vE$ tendría dos direcciones en un punto y no estaría unívocamente definido.
  \item Salen de las cargas positivas y terminan en las negativas (son líneas abiertas), y se dibujan de modo que su densidad sea
        proporcional a $\norm{\vE}$: el campo es más intenso donde están más juntas.
  \item Las superficies equipotenciales son perpendiculares en cada punto a las líneas de fuerza.
  \item Las líneas de fuerza apuntan siempre hacia donde el potencial disminuye.
\end{enumerate}
\end{proposicion}
\begin{proof}
(c) De $\phi(\vc{r})=-\int\vE\cdot\dif\vc{r}+\text{cte}$ se tiene $\dif\phi=-\vE\cdot\dif\vc{r}$. Si nos movemos sobre una superficie equipotencial,
$\dif\phi=0$, luego $\vE\perp\dif\vc{r}$ para todo desplazamiento tangente a la superficie. (d) Si nos desplazamos en el sentido del campo,
$\dif\vc{r}\parallel\vE$ y $\dif\phi=-\norm{\vE}\,\dif r<0$.
\end{proof}

La \cref{fig:lineas} muestra las líneas de campo y las equipotenciales de una carga puntual, de un dipolo y de dos cargas iguales.

% Figura: líneas de campo y superficies equipotenciales (corte meridiano).
\begin{figure}[ht]
\centering
\IfFileExists{rapido.flag}{\fbox{\parbox{0.9\textwidth}{\centering\small Figura omitida (modo rápido)}}}{%
\newcommand{\marcoL}{\draw[grispizarra, line width=0.5pt] (-3,-3) rectangle (3,3);}
\tikzset{lcampo/.style={cm1, line width=0.55pt, postaction={decorate}, decoration={markings, mark=at position 0.5 with {\arrow{Stealth[length=3.6pt]}}}},
         lequi/.style={azulacero, line width=0.5pt}}
\begin{minipage}[t]{0.32\textwidth}\centering
\begin{tikzpicture}[x=0.8cm, y=0.8cm, baseline=(current bounding box.north)]
  \begin{scope}
    \clip (-3,-3) rectangle (3,3);
    \foreach \r in {0.5,1.0,1.5,2.0,2.5,3.0}{\draw[lequi] (0,0) circle (\r);}
    \foreach \a in {0,22.5,...,337.5}{%
      \draw[lcampo] (\a:0.2) -- (\a:3.3);}
  \end{scope}
  \marcoL
  \node[cargapos] at (0,0) {$+$};
\end{tikzpicture}\\[2pt]{\small\scshape (a)\; carga puntual}
\end{minipage}\hfill
\begin{minipage}[t]{0.32\textwidth}\centering
\begin{tikzpicture}[x=0.8cm, y=0.8cm, baseline=(current bounding box.north)]
  \def\qs{-1}
  \begin{scope}
    \clip (-3,-3) rectangle (3,3);
    % equipotenciales (positivas: línea continua; negativas, por simetría: discontinua)
    \foreach \xs in {-1.25,-1.4,-1.65,-2.0,-2.6}{
      \equipC{\xs}{0}{0.05}{140}{-1}
      \draw[lequi] \ruta; \begin{scope}[yscale=-1]\draw[lequi] \ruta;\end{scope}
      \begin{scope}[xscale=-1]\draw[lequi, densely dashed] \ruta;\end{scope}
      \begin{scope}[xscale=-1,yscale=-1]\draw[lequi, densely dashed] \ruta;\end{scope}}
    \draw[lequi] (0,-3) -- (0,3);
    % líneas de campo (8 conos de igual flujo desde la carga positiva)
    \foreach \k in {1,...,8}{
      \pgfmathsetmacro{\psi}{acos(1-(2*\k-1)/8)}
      \pgfmathsetmacro{\sx}{-1+0.13*cos(\psi)}\pgfmathsetmacro{\sy}{0.13*sin(\psi)}
      \lineaC{\sx}{\sy}{0.05}{200}{1}
      \draw[lcampo] \ruta;
      \begin{scope}[yscale=-1]\draw[lcampo] \ruta;\end{scope}}
  \end{scope}
  \marcoL
  \node[cargapos] at (-1,0) {$+$};
  \node[carganeg] at (1,0) {$-$};
\end{tikzpicture}\\[2pt]{\small\scshape (b)\; dipolo}
\end{minipage}\hfill
\begin{minipage}[t]{0.32\textwidth}\centering
\begin{tikzpicture}[x=0.8cm, y=0.8cm, baseline=(current bounding box.north)]
  \def\qs{1}
  \begin{scope}
    \clip (-3,-3) rectangle (3,3);
    \foreach \xs in {-1.2,-1.3,-1.45,-1.8,-2.2,-2.7}{
      \equipC{\xs}{0}{0.05}{150}{-1}
      \foreach \sxx/\syy in {1/1,1/-1,-1/1,-1/-1}{\begin{scope}[xscale=\sxx,yscale=\syy]\draw[lequi] \ruta;\end{scope}}}
    \foreach \k in {1,...,8}{
      \pgfmathsetmacro{\psi}{acos(1-(2*\k-1)/8)}
      \pgfmathsetmacro{\sx}{-1+0.13*cos(\psi)}\pgfmathsetmacro{\sy}{0.13*sin(\psi)}
      \lineaC{\sx}{\sy}{0.05}{200}{1}
      \foreach \sxx/\syy in {1/1,1/-1,-1/1,-1/-1}{\begin{scope}[xscale=\sxx,yscale=\syy]\draw[lcampo] \ruta;\end{scope}}}
  \end{scope}
  \marcoL
  \node[cargapos] at (-1,0) {$+$};
  \node[cargapos] at (1,0) {$+$};
\end{tikzpicture}\\[2pt]{\small\scshape (c)\; dos cargas iguales}
\end{minipage}
}%
\caption{Líneas de campo (púrpura, con el sentido del campo) y superficies equipotenciales (azul, vistas en
un plano meridiano) de (a)~una carga puntual, (b)~un dipolo y (c)~dos cargas iguales. Las líneas se trazan con
igual flujo entre dos consecutivas, de modo que su densidad es proporcional a $\norm{\vE}$. Salen de las cargas
positivas y llegan a las negativas, nunca se cortan y son perpendiculares a las equipotenciales, hacia donde
el potencial disminuye. En (b) las equipotenciales de potencial negativo se trazan a trazos y la mediatriz es la
de potencial nulo; en (c) hay un punto de campo nulo en el punto medio.}
\label{fig:lineas}
\end{figure}

